\section{Local approximation and open systems}\label{sec:approx}

\motiv{\vet{The bounds of the previous chapters control the commutator of an evolved observable
$A(t)$ with every observable of a distant region. Physically one wants a stronger-looking
statement: that $A(t)$ \emph{is}, up to a small error, an observable of a bounded region, one
that grows at the velocity of the bound. The two statements are equivalent up to constants, and
the device that converts one into the other is a conditional expectation $P_X$ onto the local algebra of a
region: a linear map that projects every observable onto its part living in $X$. This chapter
builds $P_X$ on the qudit lattice as a finite average over unitaries of the complement,
\vetnew{identifies it, in the tensor-product picture, with the normalized partial trace}, and proves that the distance from $M$ to $\cA(X)$ is
controlled by the commutators of $M$ with the complement, with constants $1$ and $2$. The
evolved observable is then within a light-cone tail of a local one, and the dynamics is an
\emph{approximate cellular automaton}: it maps the observables of $O$ into those of a
neighbourhood of $O$, up to an error.

The second theme is dissipation. A quantum system coupled to an environment evolves, in the
Heisenberg picture, under a Lindblad generator: a Hamiltonian part and jump terms. Such a
generator built from local pieces satisfies the hypotheses of the path-sum bound, with the
contraction of its sub-families, a consequence of complete positivity, carried as the one
hypothesis not reproved here: locality in the argument comes from a one-line commutator
identity, and the conditional expectation gives a second route to it with the norm of the term
as strength. The bound and its closed forms for $t\ge0$ therefore hold for open systems with no
new analysis.}}

\subsection{Nets with local approximation}

\vet{Throughout, $\cA$ is a unital C*-algebra, nontrivial where a unitary has to have norm one,
with a net $Z\mapsto\cA(Z)$ over the subsets of a set of sites $\Lambda$, causal where stated
(Definition~\ref{def:net}), and $N_S(c)$ is the commutator seminorm of Chapter~\ref{sec:objects}.}

\begin{definition}[local approximation]\label{def:localapprox}
\vet{The net has \emph{local approximation} if every region $X$ carries a map
$P_X:\cA\to\cA$ with
\begin{equation}\label{eq:localapprox}
  P_Xc\in\cA(X)\qquad\text{and}\qquad\norm{c-P_Xc}\le N_{X^c}(c)\qquad(c\in\cA).
\end{equation}
A map $\mathcal T:\cA\to\cA$ is \emph{$\varepsilon$-local with the neighbourhood map
$\mathcal N$} (a map from regions to regions) if for every region $O$ and every \vetnew{$c\in\cA(O)$
there is $c'\in\cA(\mathcal N(O))$ with $\norm{\mathcal Tc-c'}\le\varepsilon\norm c$}.}
\end{definition}

\vet{The inequality in \eqref{eq:localapprox} says that an element commuting with
everything outside $X$ up to $\varepsilon$ is within $\varepsilon$ of an element of $\cA(X)$: a
statement about commutators becomes a statement about distance to a subalgebra. At $X=\Lambda$
it forces $P_\Lambda c=c$ whenever $\cA(\emptyset)$ consists of scalars, so a net with local
approximation has $\cA(\Lambda)=\cA$, \vetnew{as the qudit net does}. The converse holds on every
causal net, with the factor $2$ of the commutator.}

\needspace{12\baselineskip}
\begin{proposition}[the converse estimate; the exact case]\label{prop:converse}
\vet{\begin{enumerate}[label=(\alph*),nosep]
\item On a causal net, if \vetnew{$c'\in\cA(X)$} and $Y$ is disjoint from $X$, then
  \vetnew{$N_Y(c)\le2\norm{c-c'}$} for every $c$. On a causal net with local approximation,
  $N_{X^c}(c)\le2\norm{c-P_Xc}$.
\item $\mathcal T$ is $0$-local with $\mathcal N$ if and only if
  $\mathcal T(\cA(O))\subseteq\cA(\mathcal N(O))$ for every $O$; this is the locality condition
  of a quantum cellular automaton with the neighbourhood $\mathcal N(O)$.
\end{enumerate}}
\end{proposition}

\begin{proof}
\vet{\begin{enumerate}[label=(\alph*),nosep]
\item For $b\in\cA(Y)$ with $\norm b\le1$, \vetnew{$c'$} and $b$ commute by causality, so
  \vetnew{$\comm cb=\comm{c-c'}b$ and $\norm{\comm cb}\le2\norm{c-c'}$}; take the supremum over $b$.
\item At $\varepsilon=0$ the condition reads \vetnew{$\mathcal Tc=c'$ for some
  $c'\in\cA(\mathcal N(O))$}.\qedhere
\end{enumerate}}
\end{proof}

\subsection{The conditional expectation on the qudit lattice}

\vet{Let $\Lambda$ be finite with finite nonempty alphabets $\Omega_u$, and $\Mat_\Lambda$ the qudit
algebra of Chapter~\ref{sec:qudit}. For a region $X$ write $\Omega_{X^c}$ for the letters on the
complement and, for $\eta\in\Omega_{X^c}$ and a configuration $s$, $(s|_X,\eta)$ for the
configuration with the letters of $s$ on $X$ and $\eta$ off $X$.}

\begin{definition}[conditional expectation]\label{def:condexp}
\vet{The \emph{conditional expectation} onto $X$ is the linear map $P_X$ of $\Mat_\Lambda$ with
\begin{equation}\label{eq:condexp}
  (P_XM)_{ss'}:=\ind{s\sim_Xs'}\;\frac1{|\Omega_{X^c}|}\sum_{\eta\in\Omega_{X^c}}
  M_{(s|_X,\eta),\,(s'|_X,\eta)} .
\end{equation}}
\end{definition}

\vet{In the tensor-product picture this is the normalized partial trace over the complement,
re-embedded as $|\Omega_{X^c}|^{-1}\tr_{X^c}M\otimes1$: the entry at $(s,s')$ averages the entries of
$M$ between the pairs that agree with $s$, $s'$ on $X$ and carry the same letters off $X$. The formula makes sense for any finite alphabets.}

\begin{proposition}[properties of the conditional expectation]\label{prop:condexp}
\vet{\begin{enumerate}[label=(\alph*),nosep]
\item $P_XM$ is local to $X$, and $P_XM=M$ when $M$ is local to $X$; so $P_X\circ P_X=P_X$.
\item $P_X$ is linear, $P_X1=1$ and $(P_XM)^*=P_X(M^*)$.
\item $P_\Lambda=\mathrm{id}$ and $P_\emptyset M=\tfrac{\tr M}{|\Omega_\Lambda|}\,1$.
\end{enumerate}}
\end{proposition}

\begin{proof}
\vet{\begin{enumerate}[label=(\alph*),nosep]
\item The entry \eqref{eq:condexp} vanishes unless $s\sim_Xs'$ and depends on $s$, $s'$ through
  $s|_X$, $s'|_X$: the two conditions of Definition~\ref{def:localmatrix}. If $M$ is local to
  $X$ and $s\sim_Xs'$, every term of the sum equals $M_{ss'}$, because
  $(s|_X,\eta)\sim_X(s'|_X,\eta)$ with the same restrictions to $X$ as $s$, $s'$; the average
  of $|\Omega_{X^c}|$ copies of $M_{ss'}$ is $M_{ss'}$. If $s\not\sim_Xs'$ both sides vanish.
\item Linearity is termwise; $1$ is local to every region; the conjugate of the average is the
  average of the conjugates, and $s\sim_Xs'$ is symmetric.
\item For $X=\Lambda$ there is one letter $\eta$ off $X$, the empty one, and
  $(s|_\Lambda,\eta)=s$. For $X=\emptyset$ the condition $s\sim_\emptyset s'$ is $s=s'$ and
  $\eta$ runs over all configurations.\qedhere
\end{enumerate}}
\end{proof}

\vet{Two properties remain: that $P_X$ is a contraction, and the estimate \eqref{eq:localapprox}.
Both come from a second expression of $P_X$, as an average of unitary conjugations, which needs
a group structure on the alphabets. From here on each alphabet is cyclic, $\Omega_u=\Z_{q_u}$
with $q_u\ge1$ (every finite alphabet is one up to a renaming of its letters), and letters are
added modulo $q_u$ at each site. Write $e_n(k):=e^{2\pi\ii k/n}$ for $k\in\Z_n$; it is a
character, $e_n(k+k')=e_n(k)e_n(k')$, with $e_n(k)^n=1$ and
$\overline{e_n(k)}=e_n(-k)$.}

\needspace{12\baselineskip}
\begin{proposition}[characters, clock and shift]\label{prop:clockshift}
\vet{For $n\ge1$ and $k\in\Z_n$,
\begin{equation}\label{eq:charsum}
  \sum_{l\in\Z_n}e_n(lk)=n\,\ind{k=0}.
\end{equation}
The \emph{clock} $\clock$ and the \emph{shift} $\shift$ are the $n\times n$ matrices
$\clock_{ll'}=e_n(l)\ind{l=l'}$, $\shift_{ll'}=\ind{l=l'+1}$; \vetnew{they satisfy
$\shift\shift^*=1$} and $\clock\shift=e_n(1)\,\shift\clock$.}
\end{proposition}

\begin{proof}
\vet{If $k=0$ every term of \eqref{eq:charsum} is $1$. Otherwise $e_n(k)\ne1$ is an $n$-th root
of unity, so the geometric sum $\sum_{l<n}e_n(k)^l=(e_n(k)^n-1)/(e_n(k)-1)$ vanishes.
\vetnew{The shift is a permutation matrix, so $\shift\shift^*=1$}, and
$(\clock\shift)_{ll'}=e_n(l)\ind{l=l'+1}=e_n(1)e_n(l')\ind{l=l'+1}=e_n(1)(\shift\clock)_{ll'}$.}
\end{proof}

\vetnew{Both matrices are unitary: the clock is diagonal with unimodular entries, and the
shift, a permutation matrix, satisfies $\shift^*\shift=1$ as well.}

\vet{On a region $Y$ the products of clock and shift powers over the sites of $Y$ form the
\emph{Weyl unitaries}; we define them directly by their entries, which is what every proof
uses (Figure~\ref{fig:twirl}). \vetnew{The average of the conjugations by the Weyl unitaries of
the complement of $X$ is a \emph{twirl}, the standard name in quantum information for an average
of conjugations over a group of unitaries; it will be $P_X$.} For letters $\xi\in\Omega_Y$ and a configuration $s$, $s+\xi$
is the configuration with $s(u)+\xi(u)$ on $Y$ and $s(u)$ off $Y$, and $s-\xi$ likewise; $s-s'$
is the sitewise difference of two configurations, of which the phase $\psi_\eta$ sees the
letters on $Y$ only.}

\begin{figure}[t]
\centering
% Figure: the twirl onto a region X = {2, 3} of a chain of six sites. Transcription (for cheap
% agents): six sites 0..5 on a line; the region X = {2, 3} is a dashed box; each site u outside
% X carries a small box labelled "X^{xi_u} Z^{eta_u}" (a shift power and a clock power, the
% letters xi_u, eta_u of that site), the sites inside X carry nothing. An arrow from the row of
% boxes to the formula below reads "average over all letters (xi, eta) on X^c". The formula:
% P_X M = |L|^{-2} sum_{xi,eta} V_{xi,eta} M V_{xi,eta}^*, and underneath "= the partial trace over
% X^c, re-embedded: the phases kill the entries that differ off X, the shifts average the rest".
% All labels are set in \small type.
\begin{tikzpicture}[x=1.5cm,y=1cm,
  site/.style={circle,draw,fill=white,inner sep=1.6pt,font=\small},
  wbox/.style={draw,rectangle,inner sep=2pt,minimum height=0.5cm,font=\small,fill=Blue!15},
  region/.style={rounded corners=4pt,draw,dashed,fill=none},
  lab/.style={font=\small}]
  \foreach \u in {0,1,4,5} { \node[wbox] at (\u,0.6) {$\shift^{\xi_\u}\clock^{\eta_\u}$}; }
  \foreach \u in {0,...,5} { \node[site] at (\u,0) {\u}; }
  \draw[region] (1.55,-0.4) rectangle (3.45,0.95);
  \node[lab] at (2.5,-0.65) {$X$};
  \node[lab] at (2.5,0.6) {$1$};
  \draw[->,thick] (2.5,-0.85) -- (2.5,-1.3);
  \node[lab,anchor=west] at (-0.4,-1.95) {$P_XM=\dfrac1{|\Omega_{X^c}|^2}\displaystyle\sum_{\xi,\eta}
    V_{\xi,\eta}\,M\,V_{\xi,\eta}^*$,\quad $V_{\xi,\eta}=\prod_{u\notin X}\shift^{\xi_u}\clock^{\eta_u}$};
  \node[lab,anchor=west] at (-0.4,-2.7) {phases: kill the entries with $s\not\sim_Xs'$};
  \node[lab,anchor=west] at (-0.4,-3.15) {shifts: average the rest over the letters off $X$};
\end{tikzpicture}

\caption{\vet{The twirl onto $X=\{2,3\}$ on a chain of six sites. Each site off $X$ carries a
shift power and a clock power, with the letters $\xi_u$, $\eta_u$ of that site; the sites of $X$
carry the identity. Averaging $V_{\xi,\eta}MV_{\xi,\eta}^*$ over all letters off $X$ gives the
conditional expectation \eqref{eq:condexp}: the clock phases cancel every entry between
configurations that differ off $X$, and the shifts average the remaining entries over the
letters off $X$ (Theorem~\ref{thm:twirl}).}}
\label{fig:twirl}
\end{figure}

\begin{definition}[Weyl unitaries, twirl]\label{def:weyl}
\vet{For a region $Y$ and letters $\xi,\eta\in\Omega_Y$, the \emph{Weyl unitary} $V_{\xi,\eta}$ is
the matrix
\begin{equation}\label{eq:weyl}
  (V_{\xi,\eta})_{ss'}:=\ind{s=s'+\xi}\;\psi_\eta(s'),\qquad
  \psi_\eta(s):=\prod_{u\in Y}e_{q_u}\bigl(\eta(u)\,s(u)\bigr),
\end{equation}
the shift by $\xi$ applied after the clock phases $\psi_\eta$. The \emph{twirl} onto $X$ is the
average over the Weyl unitaries of the complement $Y=X^c$,
\begin{equation}\label{eq:twirl}
  \frac1{|\Omega_{Y}|^2}\sum_{\xi\in\Omega_Y}\sum_{\eta\in\Omega_Y}V_{\xi,\eta}\,M\,V_{\xi,\eta}^* .
\end{equation}}
\end{definition}

\begin{proposition}[Weyl unitaries]\label{prop:weyl}
\vet{$V_{\xi,\eta}$ is local to $Y$ and unitary, and for every $M$ and configurations $s,s'$
\begin{equation}\label{eq:weylconj}
  (V_{\xi,\eta}MV_{\xi,\eta}^*)_{ss'}=\psi_\eta(s-s')\;M_{s-\xi,\,s'-\xi} .
\end{equation}}
\end{proposition}

\begin{proof}
\vet{The phase $\psi_\eta(s)$ depends on $s|_Y$ only and is multiplicative in $s$, with
$\overline{\psi_\eta(s)}=\psi_\eta(-s)$ and $|\psi_\eta|=1$. The entry \eqref{eq:weyl} vanishes
unless $s\sim_Ys'$, since $s'+\xi\sim_Ys'$, and depends on the restrictions to $Y$ alone; so
$V_{\xi,\eta}$ is local to $Y$. In $V^*V$ and $VV^*$ the sums over the middle configuration
have one term, at $s+\xi$ and at $s-\xi$ respectively, where the phase meets its conjugate; both
products are the identity. In $(VMV^*)_{ss'}=\sum_{r,r'}V_{sr}M_{rr'}\overline{V_{s'r'}}$ the
only term is $r=s-\xi$, $r'=s'-\xi$, with the phase
$\psi_\eta(s-\xi)\overline{\psi_\eta(s'-\xi)}=\psi_\eta(s-s')$, the $\xi$ cancelling.}
\end{proof}

\needspace{10\baselineskip}
\begin{theorem}[the twirl is the conditional expectation]\label{thm:twirl}
\vet{For cyclic alphabets the twirl \eqref{eq:twirl} onto $X$ equals $P_XM$ for every $M$.}
\end{theorem}

\begin{proof}
\vet{By \eqref{eq:weylconj} the entry of the twirl at $(s,s')$ is
\[
  \frac1{|\Omega_Y|^2}\sum_{\xi}M_{s-\xi,s'-\xi}\sum_{\eta}\psi_\eta(s-s'),\qquad\text{and}\qquad
  \sum_\eta\psi_\eta(s-s')=\prod_{u\in Y}\;\sum_{l\in\Z_{q_u}}e_{q_u}\bigl(l\,(s(u)-s'(u))\bigr)
\]
equals $|\Omega_Y|\,\ind{s\sim_Xs'}$ by \eqref{eq:charsum} at each site. For $s\sim_Xs'$ the letters of $s$ and $s'$ off $X$ agree,
so $s-\xi=(s|_X,\zeta)$ and $s'-\xi=(s'|_X,\zeta)$ with $\zeta:=s|_Y-\xi$, and as $\xi$ runs over
$\Omega_Y$ so does $\zeta$. The remaining sum is that of \eqref{eq:condexp}, with the
normalization $|\Omega_Y|/|\Omega_Y|^2=1/|\Omega_Y|$.}
\end{proof}

\needspace{10\baselineskip}
\begin{proposition}[contraction and bimodule property]\label{prop:bimodule}
\vet{For cyclic alphabets, $\norm{P_XM}\le\norm M$, and $P_X(aMb)=a\,(P_XM)\,b$ for
$a,b\in\cA(X)$.}
\end{proposition}

\begin{proof}
\vet{Each term $V_{\xi,\eta}MV_{\xi,\eta}^*$ of the twirl has the norm of $M$, and the norm of an
average is at most the average of the norms. For $a,b\in\cA(X)$ and $V$ local to $X^c$,
$V(aMb)V^*=a(VMV^*)b$ by causality (Theorem~\ref{thm:quditnet}), termwise in the average.}
\end{proof}

\begin{theorem}[local approximation]\label{thm:localapprox}
\vet{For cyclic alphabets, every $M\in\Mat_\Lambda$ and every region $X$,
\begin{equation}\label{eq:localapproxqudit}
  \norm{M-P_XM}\le N_{X^c}(M)\le2\,\norm{M-P_XM}.
\end{equation}
The qudit net with cyclic alphabets has local approximation with $P_X$.}
\end{theorem}

\begin{proof}
\vet{Since $VV^*=1$, $M-VMV^*=\comm MV\,V^*$ for each Weyl unitary $V$ of the complement, so
$M-P_XM$ is the average of the $\comm MV\,V^*$. Each has norm at most
$\norm{\comm MV}\le N_{X^c}(M)$, because $V\in\cA(X^c)$ has norm one; so does the average. The
second inequality is Proposition~\ref{prop:converse}(a) with \vetnew{$c=M$ and
$c'=P_XM\in\cA(X)$}.}
\end{proof}

\vet{The constant $2$ is attained: for a single qubit outside $X$ and $M=\sigma^x$ at that site,
$P_XM=0$, $\norm{M-P_XM}=1$, and $N_{X^c}(M)=2$ (at least $\norm{\comm{\sigma^x}{\sigma^z}}=2$,
at most $2\norm M$). The constant $1$ is approached along the following family:
for $X=\emptyset$, \vetnew{one site with an alphabet of $n\ge2$ letters} and $M$ a rank-one
projection $|e\rangle\langle e|$, \vetnew{$\norm{M-P_\emptyset M}=1-1/n$} while $N_\Lambda(M)=1$ (the commutator of a projection with an
element of norm at most one has two off-diagonal blocks of norm at most one, and
$|e\rangle\langle f|+|f\rangle\langle e|$ with $f\perp e$ gives one).}

\subsection{Truncated dynamics and the approximate automaton}

\vet{Let $(S_i,L_i,\kappa_i)_{i\in I}$ be a family of local generators on a causal net with local
approximation, $A(t)=e^{tL}a$ for $a\in\cA(X)$, and let $X'\supseteq X$ be a larger region. The
walk bound of Chapter~\ref{sec:closed} controls the commutator of $A(t)$ with every observable
of $Y:=X'^c$; local approximation turns this into a distance.}

\needspace{14\baselineskip}
\begin{theorem}[truncated dynamics]\label{thm:truncated}
\vet{For $a\in\cA(X)$, $X\subseteq X'$, $Y=X'^c$ and $t\ge0$, with \vetnew{the walk terms $\omega_k(t)$ of Chapter~\ref{sec:closed} for the pair $(Y,X)$,
eq.~\eqref{eq:walkterms}},
\begin{equation}\label{eq:truncated}
  \norm{A(t)-P_{X'}A(t)}\le N_Y(A(t))\le2\norm a\sum_{k\ge0}\omega_k(t).
\end{equation}
Under bounded geometry $(g,\ell,p,\rho)$ with distance $d$ and $v=2g\ell p$, if
$(r-1)\rho<d(y,x)$ for all $y\in\bd Y$ and $x\in X\cap T$ (an integer $r\ge1$), the right side is
at most
$2\norm a\,|\bd Y|\,(vt)^re^{vt}/r!$ and, for every real $\mu$, at most
$2\norm a\,|\bd Y|\,e^{vt}e^{-\mu r}(e^{e^\mu vt}-1)$.}
\end{theorem}

\begin{proof}
\vet{The first inequality is local approximation. For the second, $N_Y(A(t))$ is the supremum
of $\norm{\comm{A(t)}b}$ over $b\in\cA(Y)$ with $\norm b\le1$, and the walk bound
(Theorem~\ref{thm:walkseries}) gives $2\norm a\norm b\,(\ind{X\cap Y\ne\emptyset}+\sum_k\omega_k(t))$
with $X\cap Y=\emptyset$. The geometric forms are those of Chapter~\ref{sec:closed} for the
walk series.}
\end{proof}

\vet{So $P_{X'}A(t)$ is an observable of $X'$ within a light-cone tail of $A(t)$. Read as a
property of the map $e^{tL}$, this is approximate locality with a neighbourhood map. The
natural map is the $r$-neighbourhood $O^{+r}:=\{w:\,d(w,x)<r\rho\text{ for some }x\in O\}$,
measured from the site into the region (the distance need not be symmetric).}

\needspace{10\baselineskip}
\begin{corollary}[approximate automaton]\label{cor:automaton}
\vet{\begin{enumerate}[label=(\alph*),nosep]
\item For $\Lambda$ finite and a neighbourhood map $\mathcal N$ with $O\subseteq\mathcal N(O)$
  for all $O$, $e^{tL}$ is $\varepsilon(t)$-local with $\mathcal N$, where $\varepsilon(t)$ is
  the largest, over the finitely many regions $O$, of the errors $2\sum_k\omega_k(t)$ of the
  pairs $(\mathcal N(O)^c,O)$.
\item Under bounded geometry with $\rho>0$ and an integer $r\ge1$, $e^{tL}$ is
  \vetnew{$\varepsilon_r(t)$-local with $O\mapsto O^{+r}$, where $\varepsilon_r(t)$ is either of}
  \begin{align}
    \varepsilon_r(t)&=2|T|\,\frac{(vt)^r}{r!}e^{vt},\label{eq:epsr}\\
    \varepsilon_r(t)&=2|T|\,e^{vt}e^{-\mu r}\bigl(e^{e^\mu vt}-1\bigr)\qquad(\mu\in\R).\label{eq:epsrcone}
  \end{align}
\end{enumerate}}
\end{corollary}

\begin{proof}
\vet{\begin{enumerate}[label=(\alph*),nosep]
\item For $c\in\cA(O)$ take $P_{\mathcal N(O)}e^{tL}c$ as the nearby element of
  Definition~\ref{def:localapprox} and apply \eqref{eq:truncated} with $X=O$,
  $X'=\mathcal N(O)$; the maximum over $O$ makes the error uniform.
\item $O\subseteq O^{+r}$ since $d(u,u)=0<r\rho$. A site $w\notin O^{+r}$ has $d(w,u)\ge r\rho
  >(r-1)\rho$ for every $u\in O$, so the geometric form applies, and $|\bd Y|\le|T|$ because the
  boundary consists of touched sites.\qedhere
\end{enumerate}}
\end{proof}

\vetnew{At fixed $t$ the first error of (b) vanishes as $r\to\infty$, since $(vt)^r/r!\to0$, and
the second, for $\mu>0$, is at most $2|T|\,e^{(1+e^\mu)vt-\mu r}$, so exponentially small in $r$
once $r>(1+e^\mu)vt/\mu$: up to time $t$ the dynamics is a cellular automaton with the
neighbourhood $O^{+r}$, $r$ of order $vt$, up to an exponentially small error.}

\begin{example}[truncation on a chain]\label{ex:truncation}
\vet{On a chain of five qubits with $H=\sum_{u=0}^{3}\bigl(\sigma^z_u\sigma^z_{u+1}
+\tfrac12\sigma^x_u\sigma^x_{u+1}\bigr)$ and $a=\sigma^x_0$, the error
$\norm{A(t)-P_{X_r}A(t)}$ for $X_r=\{0,\dots,r\}$ against the walk bound \eqref{eq:truncated}:
at $t=0.3$ the error is $0.56$, $0.086$, $0.017$ for $r=0,1,2$, against the bounds $9.6$, $4.6$,
$1.4$; at $t=0.1$ and $r=2$ the error is $6.6\cdot10^{-4}$ against $1.6\cdot10^{-2}$. The error
decays with $r$ as the bound predicts; the bound overshoots by factors between $17$ and $82$ at
these times, and the first two bounds quoted exceed the trivial bound $2\norm a$. For
$X_r=\Lambda$ the error is zero, as $A(t)$ is then local to $X_r$.}

\vetnew{On twelve qubits with the same Hamiltonian and observable the light cone is visible
(Figure~\ref{fig:longchain}). At $t=0.25$, $0.5$, $1$, $2$ the error falls below $0.1$ from
$r=1$, $2$, $4$, $7$ on: a front moving at about $3.5$ sites per unit time. Beyond the front the
error decays by a constant factor per site, as $e^{-3.07r}$, $e^{-2.38r}$, $e^{-1.66r}$,
$e^{-0.86r}$ at the four times (least squares over $r=3,\dots,10$), the rate falling as the
front approaches.

The walk bound \eqref{eq:truncated} decays in $r$ in the same way, at the
smaller rates $e^{-2.34r}$, $e^{-1.67r}$, $e^{-1.05r}$, $e^{-0.57r}$; it is below the trivial
bound $2\norm a$ from $r=2$ on at $t=0.25$ and from $r=5$ on at $t=0.5$, and trivial for every
$r\le10$ at $t=1$ and $t=2$, where its velocity has outrun the chain. It exceeds the error by
factors from $12$ to above $10^{10}$. The bounded-geometry bound with $v=16$ is trivial
throughout. Corollary~\ref{cor:automaton} and the paragraph after it predict the exponential
decay in the radius and a front moving linearly in $t$; the velocity they give exceeds the
observed one.}
\end{example}

\begin{figure}[htb]
\centering
% Figure: the twelve-qubit truncation errors of Example ex:truncation, generated by
% checks/long_chain_figure.py from the data of checks/long_chain.py. Transcription (for
% cheap agents): axes, the radius r = 0..10 of the region X_r to the right and the error
% eps_r(t) upward on a logarithmic scale from 1e-12 to 10; one solid curve with dots per
% time t = 0.25, 0.5, 1, 2 (labelled at its right end), each falling exponentially in r
% beyond a front that moves to the right with t; the walk bound dashed, for t = 0.25 and
% 0.5 only, where it is below the trivial bound 2 (a thin horizontal line at 2).
\begin{tikzpicture}[x=0.65cm,y=1cm,>={Stealth[length=5pt]}]
  \draw[->] (0,0) -- (11.6,0) node[right,font=\small] {$r$};
  \draw[->] (0,0) -- (0,6.18) node[above,font=\small] {$\delta_r(t)$};
  \draw (0.15,0.42) -- (-0.15,0.42) node[left,font=\scriptsize] {$10^{-12}$};
  \draw (0.15,1.68) -- (-0.15,1.68) node[left,font=\scriptsize] {$10^{-9}$};
  \draw (0.15,2.94) -- (-0.15,2.94) node[left,font=\scriptsize] {$10^{-6}$};
  \draw (0.15,4.20) -- (-0.15,4.20) node[left,font=\scriptsize] {$10^{-3}$};
  \draw (0.15,5.46) -- (-0.15,5.46) node[left,font=\scriptsize] {$10^{0}$};
  \draw (0,0.06) -- (0,-0.06) node[below,font=\scriptsize] {0};
  \draw (1,0.06) -- (1,-0.06) node[below,font=\scriptsize] {1};
  \draw (2,0.06) -- (2,-0.06) node[below,font=\scriptsize] {2};
  \draw (3,0.06) -- (3,-0.06) node[below,font=\scriptsize] {3};
  \draw (4,0.06) -- (4,-0.06) node[below,font=\scriptsize] {4};
  \draw (5,0.06) -- (5,-0.06) node[below,font=\scriptsize] {5};
  \draw (6,0.06) -- (6,-0.06) node[below,font=\scriptsize] {6};
  \draw (7,0.06) -- (7,-0.06) node[below,font=\scriptsize] {7};
  \draw (8,0.06) -- (8,-0.06) node[below,font=\scriptsize] {8};
  \draw (9,0.06) -- (9,-0.06) node[below,font=\scriptsize] {9};
  \draw (10,0.06) -- (10,-0.06) node[below,font=\scriptsize] {10};
  \draw[thin,black!50] (0,5.59) -- (11.3,5.59) node[right,font=\scriptsize,black] {$2\norm a$};
  \draw[Blue,thick] (0,5.325) -- (1,4.948) -- (2,4.622) -- (3,4.117) -- (4,3.698) -- (5,3.119) -- (6,2.638) -- (7,2.006) -- (8,1.479) -- (9,0.806) -- (10,0.248);
  \fill[Blue] (0,5.325) circle (1.2pt);
  \fill[Blue] (1,4.948) circle (1.2pt);
  \fill[Blue] (2,4.622) circle (1.2pt);
  \fill[Blue] (3,4.117) circle (1.2pt);
  \fill[Blue] (4,3.698) circle (1.2pt);
  \fill[Blue] (5,3.119) circle (1.2pt);
  \fill[Blue] (6,2.638) circle (1.2pt);
  \fill[Blue] (7,2.006) circle (1.2pt);
  \fill[Blue] (8,1.479) circle (1.2pt);
  \fill[Blue] (9,0.806) circle (1.2pt);
  \fill[Blue] (10,0.248) circle (1.2pt);
  \node[Blue,anchor=west,font=\scriptsize] at (10.15,0.248) {$t=0.25$};
  \draw[Blue,dashed] (2,5.364) -- (3,5.062) -- (4,4.717) -- (5,4.339) -- (6,3.933) -- (7,3.502) -- (8,3.049) -- (9,2.577) -- (10,2.075);
  \node[Blue,anchor=west,font=\scriptsize] at (10.15,2.075) {bound, $t=0.25$};
  \draw[black,thick] (0,5.427) -- (1,5.183) -- (2,4.985) -- (3,4.610) -- (4,4.318) -- (5,3.867) -- (6,3.513) -- (7,3.009) -- (8,2.609) -- (9,2.063) -- (10,1.627);
  \fill[black] (0,5.427) circle (1.2pt);
  \fill[black] (1,5.183) circle (1.2pt);
  \fill[black] (2,4.985) circle (1.2pt);
  \fill[black] (3,4.610) circle (1.2pt);
  \fill[black] (4,4.318) circle (1.2pt);
  \fill[black] (5,3.867) circle (1.2pt);
  \fill[black] (6,3.513) circle (1.2pt);
  \fill[black] (7,3.009) circle (1.2pt);
  \fill[black] (8,2.609) circle (1.2pt);
  \fill[black] (9,2.063) circle (1.2pt);
  \fill[black] (10,1.627) circle (1.2pt);
  \node[black,anchor=west,font=\scriptsize] at (10.15,1.627) {$t=0.5$};
  \draw[black,dashed] (5,5.393) -- (6,5.110) -- (7,4.804) -- (8,4.476) -- (9,4.129) -- (10,3.740);
  \node[black,anchor=west,font=\scriptsize] at (10.15,3.740) {bound, $t=0.5$};
  \draw[Blue!60!black,thick] (0,5.451) -- (1,5.362) -- (2,5.297) -- (3,5.061) -- (4,4.902) -- (5,4.584) -- (6,4.361) -- (7,3.987) -- (8,3.716) -- (9,3.300) -- (10,2.992);
  \fill[Blue!60!black] (0,5.451) circle (1.2pt);
  \fill[Blue!60!black] (1,5.362) circle (1.2pt);
  \fill[Blue!60!black] (2,5.297) circle (1.2pt);
  \fill[Blue!60!black] (3,5.061) circle (1.2pt);
  \fill[Blue!60!black] (4,4.902) circle (1.2pt);
  \fill[Blue!60!black] (5,4.584) circle (1.2pt);
  \fill[Blue!60!black] (6,4.361) circle (1.2pt);
  \fill[Blue!60!black] (7,3.987) circle (1.2pt);
  \fill[Blue!60!black] (8,3.716) circle (1.2pt);
  \fill[Blue!60!black] (9,3.300) circle (1.2pt);
  \fill[Blue!60!black] (10,2.992) circle (1.2pt);
  \node[Blue!60!black,anchor=west,font=\scriptsize] at (10.15,2.992) {$t=1$};
  \draw[black!60,thick] (0,5.441) -- (1,5.410) -- (2,5.349) -- (3,5.326) -- (4,5.314) -- (5,5.161) -- (6,5.083) -- (7,4.856) -- (8,4.724) -- (9,4.447) -- (10,4.276);
  \fill[black!60] (0,5.441) circle (1.2pt);
  \fill[black!60] (1,5.410) circle (1.2pt);
  \fill[black!60] (2,5.349) circle (1.2pt);
  \fill[black!60] (3,5.326) circle (1.2pt);
  \fill[black!60] (4,5.314) circle (1.2pt);
  \fill[black!60] (5,5.161) circle (1.2pt);
  \fill[black!60] (6,5.083) circle (1.2pt);
  \fill[black!60] (7,4.856) circle (1.2pt);
  \fill[black!60] (8,4.724) circle (1.2pt);
  \fill[black!60] (9,4.447) circle (1.2pt);
  \fill[black!60] (10,4.276) circle (1.2pt);
  \node[black!60,anchor=west,font=\scriptsize] at (10.15,4.276) {$t=2$};
\end{tikzpicture}

\caption{\vetnew{The truncation error $\delta_r(t)=\norm{A(t)-P_{X_r}A(t)}$ on a chain of
twelve qubits, $H=\sum_u(\sigma^z_u\sigma^z_{u+1}+\tfrac12\sigma^x_u\sigma^x_{u+1})$,
$a=\sigma^x_0$, $X_r=\{0,\dots,r\}$, against $r$ on a logarithmic scale, one curve per time
$t$. Dashed: the walk bound \eqref{eq:truncated} where it is below the trivial bound $2\norm a$
(at $t=1$ and $2$ it is trivial for every $r$ shown). Beyond a front that moves with $t$ the
error decays exponentially in $r$; the bound shares the decay, with a smaller rate and a far
larger prefactor.}}
\label{fig:longchain}
\end{figure}

\subsection{Open systems}

\vet{A quantum system in contact with an environment evolves, after the environment has been
traced out and under the Markov approximation, by a semigroup $e^{tL}$ of unital completely
positive maps, whose generator has the Lindblad form. In the Heisenberg picture, for a
Hamiltonian $H$ and \vetnew{finitely many} \emph{jump operators} \vetnew{$\Gamma_j$},
\begin{equation}\label{eq:lindblad}
  La=\ii\comm Ha+\vetnew{\sum_j}\Bigl(\Gamma_j^*a\Gamma_j-\tfrac12\bigl(\Gamma_j^*\Gamma_ja
  +a\Gamma_j^*\Gamma_j\bigr)\Bigr).
\end{equation}
The jump terms describe the environment's action: a measurement-like kick $a\mapsto\Gamma^*a\Gamma$
balanced by the anticommutator that keeps $L1=0$. On a lattice the Hamiltonian and the jumps
are local, and the generator is a sum of local terms.}

\needspace{8\baselineskip}
\begin{definition}[Lindblad terms and families]\label{def:lindbladterm}
\vet{A \emph{Lindblad term} on a finite region $S$ is a self-adjoint $H\in\cA(S)$ with \vetnew{finitely many jump
operators $\Gamma_j\in\cA(S)$}; its generator $L_S$ is \eqref{eq:lindblad}, and
the \emph{dissipator} of one jump is
$\Delta_\Gamma a:=\Gamma^*a\Gamma-\tfrac12(\Gamma^*\Gamma a+a\Gamma^*\Gamma)$, so that
$L_S=\ad_H+\sum_j\Delta_{\Gamma_j}$ \vetnew{(with no jumps the term is a Hamiltonian one)}. A
\emph{Lindblad family} is a finite family $(S_i,L_i)_{i\in I}$ of Lindblad terms, with
Hamiltonian parts $H_i$ and \vetnew{jumps $\Gamma_{ij}$, finitely many $j$ for each $i$}, together with the hypothesis
that every sub-family is \emph{contractive}: $\norm{e^{tL_J}}\le1$ for all $J\subseteq I$ and
$t\ge0$.}
\end{definition}

\vet{The contraction hypothesis is physics: $e^{tL_J}$ is unital and completely positive, and
a unital positive map between C*-algebras has norm one. Neither that norm theorem nor the
product formula that passes from the terms to their sum is reproved here; the hypothesis is
carried explicitly, and one case, a single dephasing jump, is proved below
(Proposition~\ref{prop:dephasing}).}

\needspace{10\baselineskip}
\begin{proposition}[Lindblad terms]\label{prop:lindbladterm}
\vet{For a Lindblad term on $S$ and a causal net:
\begin{enumerate}[label=(\alph*),nosep]
\item $L_S$ is bounded, $\norm{L_S}\le2\norm H+2\sum_j\norm{\Gamma_j}^2$;
\item $L_Sc=0$ for $c\in\cA(S^c)$;
\item $L_Sc\in\cA(Z\cup S)$ for $c\in\cA(Z)$ (support closure, G2);
\item $(L_Sa)^*=L_S(a^*)$.
\end{enumerate}}
\end{proposition}

\begin{proof}
\vet{\begin{enumerate}[label=(\alph*),nosep]
\item $\norm{\ad_Ha}\le2\norm H\norm a$ by \eqref{eq:ad}; the three maps
  $a\mapsto\Gamma^*a\Gamma$, $\tfrac12\Gamma^*\Gamma a$, $\tfrac12a\Gamma^*\Gamma$ have norms at most
  $\norm\Gamma^2$, $\tfrac12\norm\Gamma^2$, $\tfrac12\norm\Gamma^2$.
\item $c$ commutes with $H$, $\Gamma_j$ and $\Gamma_j^*$ by causality ($\cA(S)$ is closed under
  adjoints), so $\comm Hc=0$ and
  $\Gamma_j^*c\Gamma_j=\Gamma_j^*\Gamma_jc=\tfrac12(\Gamma_j^*\Gamma_jc+c\Gamma_j^*\Gamma_j)$.
\item Every factor lies in $\cA(Z\cup S)$ by isotony, which is closed under products and sums.
\item $(\ii\comm Ha)^*=\ii\comm H{a^*}$ for self-adjoint $H$, $(\Gamma^*a\Gamma)^*=\Gamma^*a^*\Gamma$,
  and the two halves of the anticommutator exchange.\qedhere
\end{enumerate}}
\end{proof}

\begin{theorem}[locality in the argument for Lindblad terms]\label{thm:lindbladG1}
\vet{For a Lindblad term on $S$ and every $c$:
\begin{enumerate}[label=(\alph*),nosep]
\item on any net, $\norm{\Delta_\Gamma c}\le\norm\Gamma^2N_S(c)$ for $\Gamma\in\cA(S)$, hence
  \begin{equation}\label{eq:lindbladG1direct}
    \norm{L_Sc}\le\bigl(\norm H+\textstyle\sum_j\norm{\Gamma_j}^2\bigr)N_S(c);
  \end{equation}
\item on a causal net with local approximation, $L_Sc=L_S(c-P_{S^c}c)$ and
  \begin{equation}\label{eq:lindbladG1}
    \norm{L_Sc}\le\norm{L_S}\,N_S(c)\le\bigl(2\norm H+2\textstyle\sum_j\norm{\Gamma_j}^2\bigr)N_S(c).
  \end{equation}
\end{enumerate}}
\end{theorem}

\begin{proof}
\vet{\begin{enumerate}[label=(\alph*),nosep]
\item The dissipator is a sum of two commutators,
  $\Delta_\Gamma c=\tfrac12\Gamma^*\comm c\Gamma+\tfrac12\comm{\Gamma^*}c\,\Gamma$ (expand both
  sides), and $\norm{\comm cx}\le\norm x\,N_S(c)$ for $x\in\cA(S)$
  (Proposition~\ref{prop:derivation}), applied to $x=\Gamma$ and $x=\Gamma^*$; the Hamiltonian
  part is Proposition~\ref{prop:derivation} again.
\item $P_{S^c}c\in\cA(S^c)$, so $L_S(P_{S^c}c)=0$ by Proposition~\ref{prop:lindbladterm}(b) and
  $L_Sc=L_S(c-P_{S^c}c)$ by linearity. Then
  $\norm{L_Sc}\le\norm{L_S}\norm{c-P_{S^c}c}\le\norm{L_S}N_S(c)$ by \eqref{eq:localapprox} for
  the region $S^c$, whose complement is $S$.\qedhere
\end{enumerate}}
\end{proof}

\vet{Route (a) is the sharper one: it needs neither causality nor the conditional expectation,
and with no jumps it returns the Hamiltonian strength $\norm H$ of Chapter~\ref{sec:qudit}.
Route (b) is the mechanism in one line: a term acts only on what an observable does inside its
support, the conditional expectation isolates that part, and the strength is the norm of the
term, the form in which the literature states the open-system bound.}

\needspace{14\baselineskip}
\begin{theorem}[open systems as local generators]\label{thm:opensystems}
\vet{On a causal net, a Lindblad family $(S_i,L_i)_{i\in I}$ is a family of local generators
with the strengths $\kappa_i=\norm{H_i}+\sum_j\norm{\Gamma_{ij}}^2$; with local approximation
also with $\kappa_i=\norm{L_i}$, and with $\kappa_i=2\norm{H_i}+2\sum_j\norm{\Gamma_{ij}}^2$.
Consequently the path-sum bound and its closed forms for $t\ge0$
(Chapters~\ref{sec:master}--\ref{sec:closed}\vetnew{)} hold for $A(t)=e^{tL}a$: for $a\in\cA(X)$, $b\in\cA(Y)$,
\[
  \norm{\comm{A(t)}b}\le2\norm a\norm b\Bigl(\ind{X\cap Y\ne\emptyset}+\sum_{k\ge0}\omega_k(t)\Bigr),
\]
\vetnew{where the walk terms $\omega_k(t)$ of \eqref{eq:walkterms} carry the strengths used: the
direct strengths $\norm{H_i}+\sum_j\norm{\Gamma_{ij}}^2$, or, under local approximation,
$\norm{L_i}$ or its bound $2\norm{H_i}+2\sum_j\norm{\Gamma_{ij}}^2$;} the light cone under
bounded geometry, and the truncated dynamics of
Theorem~\ref{thm:truncated}.}
\end{theorem}

\begin{proof}
\vet{G1 is Theorem~\ref{thm:lindbladG1}, (a) or (b), G2 is Proposition~\ref{prop:lindbladterm}(c),
G3 is the contraction hypothesis; the largest strengths dominate $\norm{L_i}$, and G1 is
monotone in $\kappa$.}
\end{proof}

\vetnew{The bounds are for $t\ge0$. In general a Lindblad family has no contractive negation, and
the time-reversed forms of Theorem~\ref{thm:reverse} are not available: for the dephasing term of
Proposition~\ref{prop:dephasing} below and an $a$ with $\Gamma a\Gamma=-a$ (such as $\sigma^x$ for
$\Gamma=\sigma^z$), \eqref{eq:dephasing} at the time $-t<0$ gives $e^{-t\Delta_\Gamma}a=e^{2t}a$,
of norm larger than $\norm a$. A family without jumps is a Hamiltonian family, whose negation is
contractive (Theorem~\ref{thm:hamgenerators}).}

\needspace{10\baselineskip}
\begin{proposition}[a dephasing term is contractive]\label{prop:dephasing}
\vet{For a self-adjoint jump $\Gamma$ with $\Gamma^2=1$ and no Hamiltonian part,
$\Delta_\Gamma a=\Gamma a\Gamma-a$, $\Delta_\Gamma^2=-2\Delta_\Gamma$, and for every real $t$
\begin{equation}\label{eq:dephasing}
  e^{t\Delta_\Gamma}a=\tfrac{1+e^{-2t}}2\,a+\tfrac{1-e^{-2t}}2\,\Gamma a\Gamma ;
\end{equation}
for $t\ge0$ this is a convex combination of $a$ and a unitary conjugate, so
$\norm{e^{t\Delta_\Gamma}}\le1$.}
\end{proposition}

\begin{proof}
\vet{Such a $\Gamma$ is a two-valued observable, and the term is pure dephasing in its eigenbasis.
With $\Gamma^*\Gamma=1$ the dissipator is $\Theta-\mathrm{id}$ for $\Theta:a\mapsto\Gamma a\Gamma$,
and $\Theta^2=\mathrm{id}$, so $\Delta_\Gamma^{n+1}=(-2)^n\Delta_\Gamma$. Summing the exponential
series, $e^{t\Delta_\Gamma}=\mathrm{id}+\sum_{n\ge0}\frac{t^{n+1}(-2)^n}{(n+1)!}\Delta_\Gamma
=\mathrm{id}+\frac{1-e^{-2t}}2\Delta_\Gamma$, which is \eqref{eq:dephasing}. Both $\mathrm{id}$
and $\Theta$ have norm at most one ($\Gamma$ is unitary), and the coefficients are nonnegative
for $t\ge0$.}
\end{proof}

\begin{example}[two qubits with dephasing]\label{ex:dephasing}
\vet{Two qubits, $H=\sigma^x_0\sigma^x_1$, one jump $\Gamma=\sigma^z_0$, and the observables
$a=\sigma^z_0$, $b=\sigma^z_1$ of Example~\ref{ex:twoqubit}. The dissipator vanishes on $a$,
since $a$ commutes with $\Gamma$, so $La=\ii\comm Ha$ and the derivative of $\comm{A(t)}b$ at
$t=0$ is the nonzero double commutator of Chapter~\ref{sec:qudit}; hence $\comm{A(t)}b\ne0$ for
some $t>0$ (the contraction of this term, whose Hamiltonian part does not commute with the jump,
is the carried hypothesis, not Proposition~\ref{prop:dephasing}).}
\end{example}

\vetnew{For a three-qubit family with
random Hamiltonian and jump terms on $\{0,1\}$ and $\{1,2\}$ and $a=\sigma^x_0$,
$\norm{\comm{A(t)}{\sigma^x_2}}$ is $0$ at $t=0$ and $0.11$ at $t=0.5$; for the same $a$ and
$b=\sigma^z_2$ at $t=0.3$ the commutator norm is $0.026$ against the walk bound $4.5$ with the
strengths $\norm{L_i}$ and $272$ with the strengths $2\norm{H_i}+2\sum_j\norm{\Gamma_{ij}}^2$.
Both bounds exceed the trivial bound $2$; the first overshoots by a factor of $170$, more than
the Hamiltonian bounds above, and the larger strengths cost a further factor of $60$ (the
direct strengths of Theorem~\ref{thm:lindbladG1}(a) are half of them, and the series they enter
is correspondingly smaller).}

\subsection*{Notes and references}
\vetnew{\emph{Local approximation.} The review \cite{chen-lucas-yin}, its Proposition~4.1, states
for the expectation $\mathbb E_S$ onto a region $S$, the Haar average over the unitaries of the
complement of $S$,
\[
  \norm{\mathbb E_SA-A}\le\sup_B\norm{\comm AB}\le2\norm{\mathbb E_SA-A},
\]
the supremum over the unit operators $B$ on the complement. Theorem~\ref{thm:localapprox} is
this statement with $S=X$, $A=M$ and $\mathbb E_S=P_X$: the Haar average and the finite twirl
over the Weyl unitaries are both the normalized partial trace over the complement, re-embedded
(Theorem~\ref{thm:twirl}), and the supremum is $N_{X^c}(M)$.

\emph{Truncated dynamics.} The review's Proposition~4.3
bounds the distance from $\tau_t(A)$ to the evolution under the Hamiltonian restricted to a
larger region; Theorem~\ref{thm:truncated} matches it in shape, an observable of the larger
region within the tail of the walk series, with the conditional expectation $P_{X'}A(t)$ as the
local approximant and the constants of these notes.

\emph{Open systems.} The review's Theorem~3.14 is its Lieb-Robinson bound for a local Lindblad
generator, for $t\ge0$, with the norms of the Lindblad terms as strengths. On the qudit net with
cyclic alphabets, which has local approximation (Theorem~\ref{thm:localapprox}),
Theorem~\ref{thm:opensystems} with $\kappa_i=\norm{L_i}$ and Proposition~\ref{prop:finite} give
it in the notation of the review's (3.36) recalled at the end of Chapter~\ref{sec:qudit}:
\[
  \mathcal C_{AB}(t)\le\sum_{u\in A,\,v\in B}\bigl(e^{2th}\bigr)_{uv}\quad(t\ge0),\qquad
  h_{uv}=\sum_{i:\,u,v\in S_i}\norm{L_i},
\]
by the argument given there. The review obtains the contraction of the sub-families from
complete positivity; here it is the hypothesis carried in Definition~\ref{def:lindbladterm}.}

\paragraph{Checks.} {\small The twirl against the partial trace (exactly for qubits, numerically
for qutrits), the properties of the conditional expectation, the estimate
\eqref{eq:localapproxqudit} with the attained factor $2$, the truncation errors on chains of
three and five qubits \vetnew{(section~3 of the script; the chapter quotes the five-site
numbers)}, the contractivity, the vanishing off the support, the mechanism
$L_Sc=L_S(c-P_{S^c}c)$, the strengths and the walk bound for the random Lindblad family\vetnew{ (section~4 of the script)}, and the
degenerate cases ($X=\Lambda$, $X=\emptyset$, no jumps, $t=0$) are computed by
\texttt{checks/m3\_check.py}. \vetnew{The twelve-qubit errors and bounds of
Example~\ref{ex:truncation} are computed by \texttt{checks/long\_chain.py}; the fronts, decay rates, ratios and Figure~\ref{fig:longchain} by
\texttt{checks/long\_chain\_figure.py}.}}
